
\section{Ontología temporal latente}

Los procesos constructivos utilizan objetos que desaparecen: andamios, rampas, pasarelas temporales, grúas, gatos, ataguías, moldes, soportes de erección y fijaciones provisionales.

Sea \(\mathcal T\) el conjunto de estas entidades.

\begin{definition}[Ontología temporal latente mínima]
\begin{equation}
\boxed{
\mathcal T^\star
=
\arg\min_{\mathcal T}
C(\mathcal T)
}
\end{equation}
sujeto a la existencia de una historia \(H\) realizable mediante física convencional que alcance \(X_T\).
\end{definition}

La inferencia de \(\mathcal T^\star\) no debe exceder la resolución soportada por la evidencia. El estudio adversarial desarrollado más adelante demuestra que, en numerosos casos, sólo puede recuperarse una \emph{clase funcional} y no el dispositivo histórico exacto.

\section{Jerarquía de identificabilidad}

Los benchmarks sugieren cuatro niveles:

\begin{description}
\item[I0 --- orden causal grueso:] por ejemplo, cimentación antes de torre.
\item[I1 --- clase funcional:] por ejemplo, acceso aéreo temporal.
\item[I2 --- familia de mecanismo:] por ejemplo, formación incremental de cable.
\item[I3 --- implementación histórica exacta:] tipo de aparato, geometría, fabricante y fecha.
\end{description}

\begin{principle}[Resolución de la afirmación]
Ninguna afirmación debe formularse a un nivel de detalle superior al nivel de identificabilidad soportado por las modalidades disponibles.
\end{principle}

\section{Campo de fuerza residual}

Definimos

\begin{definition}[Fuerza residual]
\begin{equation}
\boxed{
\Fres
=
\mathbf F_{\mathrm{req}}
-
\sum_k
\mathbf F_{{\mathrm{known}},k}.
}
\end{equation}
\end{definition}

El problema conjunto puede escribirse como

\begin{equation}
\boxed{
\min_{H,\mathcal T,\Fres}
\left[
\int
\|\Fres\|^2\,dV\,dt
+
\lambda_T C(\mathcal T)
+
\lambda_H C(H)
\right]
}
\end{equation}

sujeto a estado terminal, equilibrio, accesibilidad, evidencia y restricciones cronológicas.

\begin{principle}[No new physics before latent process exhaustion]
Una fuerza residual no puede interpretarse como evidencia de nueva física antes de demostrar que sobrevive a:
\begin{enumerate}
\item refinamiento ontológico;
\item descomposición modular;
\item objetos temporales;
\item mecanismos convencionales;
\item chequeos de identificabilidad;
\item propagación de incertidumbre.
\end{enumerate}
\end{principle}

\section{Extensión contrafactual: compensación gravitatoria}

Este apartado define una familia de hipótesis falsables. No constituye evidencia histórica ni física establecida.

Introducimos un campo adimensional \(\alpha(\mathbf x,t)\):

\begin{equation}
\mathbf g_{\mathrm{eff}}
=
-(1-\alpha)g\hat{\mathbf z}.
\end{equation}

Los regímenes son:

\begin{align}
\alpha=0 &: \text{gravedad ordinaria},\\
0<\alpha<1 &: \text{reducción efectiva de peso},\\
\alpha=1 &: \text{neutralización ideal},\\
\alpha>1 &: \text{aceleración efectiva ascendente}.
\end{align}

La energía potencial efectiva uniforme sería

\begin{equation}
U_{\mathrm{eff}}
=
(1-\alpha)mgz.
\end{equation}

En una rampa,

\begin{equation}
F_\alpha
=
(1-\alpha)mg
(\sin\theta+\mu\cos\theta),
\end{equation}

mientras el contacto con la superficie siga siendo válido.

Si existe una fuerza convencional máxima \(F_{\max}\),

\begin{equation}
\boxed{
\alpha_{\min}
=
1-
\frac{
F_{\max}
}{
mg(\sin\theta+\mu\cos\theta)
}.
}
\end{equation}

Este parámetro puede utilizarse como \emph{control contrafactual}: si todas las historias físicamente convencionales producen residuo nulo, entonces \(\alpha_{\min}=0\) es suficiente.

\section{Extensión contrafactual: campo elástico-helicoidal}

Para formalizar una interacción restauradora tipo resorte con geometría helicoidal, definimos en coordenadas cilíndricas

\begin{equation}
\xi(r,\phi,z)
=
r-r_0-b\phi+\lambda z
\end{equation}

y

\begin{equation}
U_{\mathrm{SE}}
=
\frac12k\xi^2.
\end{equation}

La fuerza conservativa asociada es

\begin{equation}
\mathbf F_{\mathrm{SE}}
=
-\nabla U_{\mathrm{SE}}.
\end{equation}

Como

\begin{equation}
\nabla\xi
=
\hat{\mathbf e}_r
-
\frac{b}{r}\hat{\mathbf e}_\phi
+
\lambda\hat{\mathbf e}_z,
\end{equation}

se obtiene

\begin{equation}
\boxed{
\mathbf F_{\mathrm{SE}}
=
-k\xi\hat{\mathbf e}_r
+
\frac{kb\xi}{r}\hat{\mathbf e}_\phi
-
k\lambda\xi\hat{\mathbf e}_z.
}
\end{equation}

Para \(\xi<0\) y \(\lambda>0\), la componente vertical puede apuntar hacia arriba.

Este campo no debe identificarse con gravedad newtoniana. Un campo general puede descomponerse formalmente como

\begin{equation}
\Fres
=
-\nabla\Phi
+
\nabla\times\mathbf A.
\end{equation}

La componente rotacional pertenece a una clase de interacción más general.

Para una curva elástica puede utilizarse energía de flexión

\begin{equation}
E_b
=
\frac B2
\int
\kappa(s)^2ds.
\end{equation}

La familia elástico-helicoidal se conserva únicamente como modelo contrafactual a ajustar si apareciera un residuo robusto futuro. Ningún benchmark actual la requiere.
